\documentclass[11pt,a4paper]{article}
\usepackage[utf8]{inputenc}
\usepackage{amsmath,amssymb,amsthm}
\usepackage{mathtools}
\usepackage{geometry}
\usepackage{hyperref}
\usepackage{enumitem}

\geometry{margin=2.5cm}

\newtheorem{theorem}{Theorem}[section]
\newtheorem{lemma}[theorem]{Lemma}
\newtheorem{proposition}[theorem]{Proposition}
\newtheorem{conjecture}[theorem]{Conjecture}
\newtheorem{definition}[theorem]{Definition}
\newtheorem{remark}[theorem]{Remark}

\newcommand{\C}{\mathbb{C}}
\newcommand{\R}{\mathbb{R}}
\newcommand{\N}{\mathbb{N}}
\newcommand{\Z}{\mathcal{Z}}
\newcommand{\LOp}{\mathcal{L}}
\newcommand{\TOp}{T}
\newcommand{\spec}{\operatorname{spec}}
\newcommand{\supp}{\operatorname{supp}}
\newcommand{\tr}{\operatorname{tr}}
\newcommand{\Tr}{\operatorname{Tr}}
\newcommand{\RE}{\operatorname{Re}}
\newcommand{\IM}{\operatorname{Im}}

\title{\bf Objective 1 -- Progress Report:\\
Asymptotic Estimates for the Iota Walk}
\author{}
\date{\today}

\begin{document}
\maketitle

\begin{abstract}
We report on the status of Objective~1 of the research programme, which aims to establish asymptotic estimates for the partial sums (the iota walk) of the Dirichlet eta function and to exploit the greedy condition (GC) to force the real part of any zero of $\eta(s)$ to be $1/2$.  While the convergence properties of the Dirichlet series are well understood, the crucial step --- proving that GC and the convergence to zero together imply $\RE(s)=1/2$ --- remains an open problem equivalent to the Riemann Hypothesis.  We summarise what is known, give precise formulations of the missing estimates, and outline how the greedy condition might be employed to break the current impasse.
\end{abstract}

\section{The Iota Walk and the Greedy Condition}
We recall the definitions.  For $s=\sigma+it\in\C$ with $\sigma>0$, set
\[
v_n(s)=(-1)^{n-1}n^{-s},\qquad 
\mathbf{S}_N(s)=\sum_{n=1}^{N}v_n(s)=(X_N(s),Y_N(s))\in\R^2.
\]
The Dirichlet eta function is $\eta(s)=\lim_{N\to\infty}\mathbf{S}_N(s)$ whenever the limit exists.  The greedy condition (GC) for a candidate zero $s$ is
\[
\langle\mathbf{S}_{N-1}(s),\,v_N(s)\rangle \le 0\qquad\forall N\ge 1.
\]

\section{Known Convergence Properties}
The alternating series $\sum_{n=1}^\infty (-1)^{n-1}n^{-s}$ converges (conditionally) for all $\sigma>0$ and diverges for $\sigma\le 0$.  The limit is $\eta(s)$, which is analytic on $\sigma>0$ and extends meromorphically to $\C$ via its relation to $\zeta(s)$:
\[
\eta(s)=(1-2^{1-s})\zeta(s).
\]
Thus, the zeros of $\eta(s)$ on $\sigma>0$ are exactly the zeros of $\zeta(s)$ (the non‑trivial zeros) together with the pole of $\zeta$ at $s=1$ where the factor $1-2^{1-s}$ vanishes.  The non‑trivial zeros lie in the critical strip $0<\sigma<1$.  The Riemann Hypothesis asserts that they all satisfy $\sigma=1/2$.

The convergence of the series is not absolute for $0<\sigma\le 1$; the partial sums exhibit oscillations of size roughly $N^{1/2-\sigma}$ before eventually settling (if $\sigma>1/2$ the oscillations decay, if $\sigma<1/2$ they grow).  However, proving that the limit cannot be zero unless $\sigma=1/2$ is exactly the Riemann Hypothesis, and it remains unproven.

\section{Available Asymptotic Estimates}
For $\sigma>1$, the series converges absolutely and $|\mathbf{S}_N(s)-\eta(s)|=O(N^{1-\sigma})$, so the limit is non‑zero because $|\eta(s)|\ge \zeta(\sigma)-1>0$ (since $\zeta(s)$ has no zeros for $\sigma>1$).  Thus GC, if it could hold at all in this region, would be irrelevant because the walk does not tend to zero.

For $0<\sigma<1$, the situation is more delicate.  Using summation by parts and the classical bound for the exponential sum $\sum_{n\le N} n^{-it}$, one obtains
\[
\mathbf{S}_N(s) = \eta(s) + R_N(s),\qquad
|R_N(s)| \le C_\sigma N^{-\sigma} \cdot \max_{M\le N} \Bigl|\sum_{n=1}^{M} n^{-it}\Bigr|.
\]
The inner sum is bounded by $O(\sqrt{M}\log M)$ (a consequence of the Weyl–van der Corput theory).  Hence
\[
|R_N(s)| = O(N^{1/2-\sigma}\log N).
\]
This shows that for $\sigma>1/2$, the remainder decays, albeit slowly.  For $\sigma<1/2$, the bound becomes unbounded, but this does not prove divergence; the series still converges conditionally, but the partial sums can stray far from the limit before eventually returning.  In particular, the bound does not prevent the limit $\eta(s)$ from being zero.

The greedy condition imposes a very strong restriction on the evolution of the walk: it forces the squared distance $|\mathbf{S}_N|^2$ to be non‑increasing.  This is a global property that is not captured by the standard estimates above.  An analytic approach to GC would require controlling the sign of $\langle\mathbf{S}_{N-1},v_N\rangle$ for all $N$, a task of formidable difficulty.

\section{The Core Unproven Estimate}
The missing estimate can be formulated as follows.

\begin{conjecture}[Greedy condition forces critical line]\label{conj:main}
Let $s=\sigma+it$ with $\sigma>0$, $\sigma\neq 1/2$, and suppose that the iota partial sums $\mathbf{S}_N(s)$ satisfy the greedy condition for all $N\ge 1$.  Then $\lim_{N\to\infty}\mathbf{S}_N(s)\neq 0$.
\end{conjecture}

In other words, if the walk is ``greedy'' (never steps towards its current position), it cannot end at the origin unless $\sigma=1/2$.  This conjecture, if true, together with the known fact that all zeros on the critical line satisfy GC (numerically checked for the first several million zeros), would prove the Riemann Hypothesis.

Note that the converse direction is easier: if the Riemann Hypothesis holds, then all non‑trivial zeros lie on $\sigma=1/2$.  Proving that they also satisfy GC is a separate problem, but one that might be approachable via the Riemann–Siegel formula or the integral representations.

\section{Partial Results Towards Conjecture~\ref{conj:main}}
We sketch a few partial results that have been obtained.

\begin{lemma}[GC and monotonicity of the squared distance]\label{lem:monotone}
For any $s$, the greedy condition implies
\[
|\mathbf{S}_N(s)|^2 \le |\mathbf{S}_{N-1}(s)|^2 + |v_N(s)|^2,
\]
and in particular $|\mathbf{S}_N(s)|^2$ cannot increase by more than $|v_N|^2$ at each step.  If GC holds, then $|\mathbf{S}_N|^2 \le |\mathbf{S}_{N-1}|^2$ holds whenever $\langle\mathbf{S}_{N-1},v_N\rangle \le -|v_N|^2/2$, which is a weaker condition.
\end{lemma}
\begin{proof}
From $|\mathbf{S}_N|^2 = |\mathbf{S}_{N-1}|^2 + 2\langle\mathbf{S}_{N-1},v_N\rangle + |v_N|^2$, the GC $\langle\mathbf{S}_{N-1},v_N\rangle\le 0$ gives $|\mathbf{S}_N|^2 \le |\mathbf{S}_{N-1}|^2 + |v_N|^2$.  The second statement is elementary.
\end{proof}

\begin{lemma}[Divergence for small $\sigma$]\label{lem:small-sigma}
If $0<\sigma<1/2$ and GC holds, then $|\mathbf{S}_N(s)|^2$ cannot tend to zero because $|v_N|=N^{-\sigma}$ does not tend to zero fast enough to compensate for the cumulative effect of the steps, regardless of their directions.
\end{lemma}
\begin{proof}[Sketch]
Assume $|\mathbf{S}_N|\to 0$.  From Lemma~\ref{lem:monotone}, $|\mathbf{S}_N|^2 \le |\mathbf{S}_{N-1}|^2 + N^{-2\sigma}$.  Summing this inequality from $N=1$ to $M$ gives $|\mathbf{S}_M|^2 \le \sum_{n=1}^{M} n^{-2\sigma}$.  Since $\sigma<1/2$, the series $\sum n^{-2\sigma}$ diverges, so the bound is useless.  To obtain a contradiction, one needs a lower bound.  The GC does not obviously provide one.
\end{proof}

Thus, the case $\sigma<1/2$ does not yield a contradiction from GC alone.  A deeper analysis of the directions of the steps is needed.

\begin{lemma}[Conditional limit identity]\label{lem:limit}
If $\lim_{N\to\infty}\mathbf{S}_N(s)=0$, then for any sequence $N_k\to\infty$,
\[
\lim_{k\to\infty} \frac{1}{N_k}\sum_{n=1}^{N_k} \mathbf{S}_{n-1}(s)\cdot v_n(s) = -\frac12\lim_{k\to\infty}\frac{1}{N_k}\sum_{n=1}^{N_k}|v_n(s)|^2,
\]
provided the limits exist.  If additionally GC holds, then the left‑hand side is $\le 0$, and we obtain
\[
\liminf_{N\to\infty}\frac{1}{N}\sum_{n=1}^{N} |v_n(s)|^2 \ge 0,
\]
which is automatically true, giving no new restriction.
\end{lemma}
\begin{proof}
From $|\mathbf{S}_N|^2 = |\mathbf{S}_{N-1}|^2 + 2\langle\mathbf{S}_{N-1},v_N\rangle + |v_N|^2$, sum and divide by $N$, and use $|\mathbf{S}_N|\to0$ to eliminate the $|\mathbf{S}_N|^2$ terms.
\end{proof}

These rather weak consequences illustrate the difficulty: the greedy condition is a pointwise inequality, but the limit behaviour involves cancellations over large ranges.

\section{Outlook}
To make progress on Conjecture~\ref{conj:main}, we propose to investigate the following avenues:
\begin{itemize}
    \item \textbf{Exponential sum estimates with GC:} Use the specific alternating sign pattern to link GC to the lack of sign changes in the correlation between the partial sums and the steps.  This might lead to a proof that $\langle\mathbf{S}_{N-1},v_N\rangle$ must be positive infinitely often unless $\sigma=1/2$.
    \item \textbf{Continuous analogue:} The continuous greedy condition derived from Riemann's integral representation may be more tractable, as it becomes a differential inequality.  A proof that the continuous walk cannot be monotone decreasing to zero unless $\sigma=1/2$ would imply Conjecture~\ref{conj:main} via a suitable discretisation.
    \item \textbf{Connection to the GDST operator:} The correlation operator $T_\sigma$ and its Fredholm determinant provide an exact relation between the eigenvalues and $\zeta(s)$.  Perhaps the greedy condition can be translated into a positivity property of the spectral measure, which together with the determinant identity would force the zeros to be on the line.  However, the fatal error of the original GDST proof must be avoided.
\end{itemize}

\section{Conclusion}
Objective~1 remains largely unfulfilled.  The available asymptotic estimates do not suffice to prove the crucial implication that GC and convergence to zero imply $\sigma=1/2$.  The task is equivalent to the Riemann Hypothesis itself.  The greedy condition provides a new geometric angle, but its exploitation requires significant new ideas.  The above lemmas are only the first steps; a complete solution remains a major open problem.

\begin{thebibliography}{9}
\bibitem{Titchmarsh} E.\,C.\ Titchmarsh, \textit{The Theory of the Riemann Zeta‑function}, 2nd ed., Oxford Univ.\ Press, 1986.
\bibitem{Iota} V.\ Geere, \textit{On the Iota Function and the Greedy Condition: A Reformulation, Not a Proof}, manuscript, 2026.
\end{thebibliography}

\end{document}